\documentclass[12pt]{article}

\usepackage{amsmath, amssymb}
\usepackage{geometry}
\usepackage{array}
\usepackage{booktabs}
\usepackage{physics}
\usepackage{hyperref}

\geometry{margin=1in}

\title{Assignment 1}
\author{Zachary Tullier}
\date{6/29/26}

\begin{document}

\maketitle

\section*{Executive Summary}

Radiation is best understood by separating \textbf{non-ionizing radiation}, which primarily excites molecules or heats matter, from \textbf{ionizing radiation}, whose individual quanta carry enough energy to remove bound electrons from atoms or molecules. In public-health guidance, radiofrequency radiation, microwaves, infrared radiation, visible light, and commonly ultraviolet radiation are treated as non-ionizing, whereas X-rays, gamma rays, and particulate radiations such as alpha particles, beta particles, and neutrons are treated as ionizing. The exact named-band boundaries of the electromagnetic spectrum are conventional rather than immutable physical cutoffs, but X-rays consistently occupy the short-wavelength, high-frequency, high-energy region.

For electromagnetic radiation, wavelength $\lambda$, frequency $f$, and photon energy $E$ are related by
\begin{equation}
    c = \lambda f,
\end{equation}
and
\begin{equation}
    E = hf = \frac{hc}{\lambda},
\end{equation}
where $c$ is the speed of light and $h$ is Planck's constant. X-rays are commonly placed at wavelengths of approximately $10^{-11}$ to $10^{-8}\,\mathrm{m}$ and photon energies from roughly $100\,\mathrm{eV}$ to $100\,\mathrm{keV}$, depending on the convention being used. Diagnostic medical X-ray tubes commonly operate with electron energies up to about $150\,\mathrm{keV}$, producing clinically useful beams mainly in the tens of keV after filtration.

For X-rays in matter, the three canonical photon interactions are \textbf{photoelectric absorption}, \textbf{Compton scattering}, and \textbf{pair production}. In diagnostic imaging, only the first two are important, because pair production requires at least $1.022\,\mathrm{MeV}$, far above ordinary diagnostic tube potentials. Photoelectric absorption is strongly favored at lower photon energies and in higher-$Z$ materials, which is why bone and high-$Z$ contrast materials attenuate diagnostic X-rays strongly. Compton scattering becomes more important as energy rises through the diagnostic range and is a major source of image-degrading scatter.

X-rays in conventional tubes are generated by two mechanisms: \textbf{bremsstrahlung}, a continuous spectrum produced when fast electrons are decelerated in the target field, and \textbf{characteristic radiation}, a discrete line spectrum produced when target-shell vacancies are filled by outer-shell electrons. The escaped beam is then shaped by tube voltage, tube current--time product, target material, anode angle, added filtration, and generator waveform. Modern diagnostic systems overwhelmingly use \textbf{rotating-anode tubes} because the rotating track spreads heat and permits higher loading than stationary-anode designs.

Biologically, ionizing radiation acts chiefly through \textbf{direct DNA damage} and \textbf{indirect damage via radiolysis of water} and reactive species. High doses can cause threshold tissue reactions such as skin injury or acute radiation syndrome; lower doses increase long-term stochastic risks such as cancer. Protection therefore centers on \textbf{ALARA}, meaning ``as low as reasonably achievable,'' and the classic triad of \textbf{time, distance, and shielding}, supported in medical imaging by beam filtration, collimation, equipment interlocks, and quality assurance.

\section{Radiation Classes and Physical Properties}

The most important conceptual distinction is not simply ``wave versus particle,'' but whether the radiation can ionize matter. Ionizing radiation removes electrons from atoms or molecules. Non-ionizing radiation has lower photon energy and usually causes excitation, vibration, rotation, or heating rather than ionization.

Electromagnetic radiation behaves as both a wave and a collection of photons. Wave properties include wavelength, frequency, speed, reflection, refraction, diffraction, interference, and polarization. Photon properties become especially important at high frequencies, where individual photon energies are large enough to cause chemical or biological damage.

\begin{longtable}{p{0.16\textwidth} p{0.16\textwidth} p{0.22\textwidth} p{0.24\textwidth} p{0.14\textwidth}}
\caption{Major radiation classes and their properties.}\\
\toprule
\textbf{Radiation type} & \textbf{Ionizing status} & \textbf{Approximate physical range} & \textbf{Dominant physical action in matter} & \textbf{Typical use or note}\\
\midrule
\endfirsthead

\toprule
\textbf{Radiation type} & \textbf{Ionizing status} & \textbf{Approximate physical range} & \textbf{Dominant physical action in matter} & \textbf{Typical use or note}\\
\midrule
\endhead

Radio waves & Non-ionizing & $\lambda > 10^{-1}\,\mathrm{m}$; $f < 3\times 10^9\,\mathrm{Hz}$ & Induced currents and field coupling & Communications and RF systems\\

Microwaves & Non-ionizing & $10^{-3}$ to $10^{-1}\,\mathrm{m}$ & Dielectric heating & Radar, ovens, telecom\\

Infrared & Non-ionizing & $7\times 10^{-7}$ to $10^{-3}\,\mathrm{m}$ & Vibrational excitation and heating & Thermal radiation\\

Visible light & Non-ionizing & $4\times 10^{-7}$ to $7\times 10^{-7}\,\mathrm{m}$ & Electronic excitation & Vision and optics\\

Ultraviolet & Usually treated as non-ionizing in public-health guidance, but spans the transition region & $10^{-8}$ to $4\times 10^{-7}\,\mathrm{m}$ & Photochemical injury; high-energy UV approaches ionizing behavior depending on context & Skin and eye effects\\

X-rays & Ionizing & About $10^{-11}$ to $10^{-8}\,\mathrm{m}$; often about $100\,\mathrm{eV}$ to $100\,\mathrm{keV}$ & Photoelectric absorption, Compton scattering, and pair production above threshold & Imaging, diffraction, therapy planning\\

Gamma rays & Ionizing & Often $\lambda < 10^{-11}\,\mathrm{m}$; $E > 100\,\mathrm{keV}$ by one common convention & Same photon interactions as X-rays, usually at higher energy & Nuclear and astrophysical sources\\

Alpha, beta, and neutron radiation & Ionizing & Usually described by particle energy rather than EM wavelength & Dense ionization tracks for charged particles; nuclear interactions for neutrons & Nuclear medicine, reactors, radioactive decay\\
\bottomrule
\end{longtable}

\section{X-Ray Interactions with Matter}

When photons interact with matter, the principal macroscopic quantity is attenuation. For a narrow photon beam passing through thickness $x$, transmission is described by
\begin{equation}
    N(x) = N_0 e^{-\mu x},
\end{equation}
where $N_0$ is the incident photon number, $N(x)$ is the transmitted photon number, and $\mu$ is the linear attenuation coefficient. A practical measure of beam quality is the \textbf{half-value layer} (HVL), the thickness of absorber required to reduce beam intensity by one half:
\begin{equation}
    \mathrm{HVL} = \frac{\ln 2}{\mu}.
\end{equation}

\subsection{Photoelectric Absorption}

The \textbf{photoelectric effect} occurs when a photon transfers all of its energy to a bound electron, usually in an inner shell. The emitted photoelectron has kinetic energy
\begin{equation}
    T_e = E_{\gamma} - B_e,
\end{equation}
where $E_{\gamma}$ is the photon energy and $B_e$ is the electron binding energy. The resulting vacancy is filled by an outer-shell electron, producing either characteristic X-rays or Auger electrons.

Photoelectric absorption is favored by low photon energy and high atomic number $Z$. This is why bone, iodine, barium, and lead attenuate diagnostic X-rays strongly.

\subsection{Compton Scattering}

\textbf{Compton scattering} is an inelastic interaction between an incoming photon and a weakly bound outer electron. Part of the photon energy is transferred to the electron, while a lower-energy photon is scattered in a new direction. The scattered photon wavelength shift is
\begin{equation}
    \Delta \lambda = \lambda' - \lambda
    = \frac{h}{m_e c}(1-\cos\theta),
\end{equation}
where $\theta$ is the scattering angle and $m_e$ is the electron mass. In diagnostic radiography, Compton scatter reduces image contrast and increases unwanted dose outside the primary beam.

\subsection{Pair Production}

\textbf{Pair production} occurs when a high-energy photon interacts in the electromagnetic field of a nucleus and converts into an electron--positron pair. The threshold energy is
\begin{equation}
    E_{\min} = 2m_e c^2 = 1.022\,\mathrm{MeV}.
\end{equation}
Because diagnostic X-ray tubes operate far below this threshold, pair production is not important in ordinary radiography.

\section{Biological Effects of Radiation}

Ionizing radiation can damage tissue by depositing energy at molecular scales. The most important target is DNA. Damage can occur directly, when ionization breaks chemical bonds in DNA, or indirectly, when radiation ionizes water molecules and produces reactive species that later attack DNA or other cellular structures.

Possible outcomes include correct repair, incorrect repair, cell death, mutation, or malignant transformation. High doses can produce deterministic or threshold effects such as skin burns, cataracts, sterility, or acute radiation syndrome. Lower doses are usually discussed in terms of stochastic effects, especially cancer risk, where the probability of harm increases with dose.

Radiation protection follows the ALARA principle:
\begin{enumerate}[label=\arabic*.]
    \item \textbf{Time:} reduce the time spent near radiation sources.
    \item \textbf{Distance:} increase distance from the source, since intensity often decreases approximately with the inverse square of distance for point-like sources.
    \item \textbf{Shielding:} place suitable absorbing material between the source and the person.
\end{enumerate}

\section{How X-Rays Are Produced}

In a conventional X-ray tube, electrons are released from a heated cathode by \textbf{thermionic emission}, accelerated through a high potential difference toward a metal anode, and rapidly decelerated in the target. The tube voltage determines the maximum possible photon energy:
\begin{equation}
    E_{\max} = eV,
\end{equation}
where $e$ is the elementary charge and $V$ is the tube potential. If the tube potential is expressed in kilovolts, the endpoint energy is approximately the same numerical value in keV. For example, a $100\,\mathrm{kVp}$ tube has a maximum photon energy of about $100\,\mathrm{keV}$.

\subsection{Bremsstrahlung Radiation}

\textbf{Bremsstrahlung}, or braking radiation, is produced when fast electrons are decelerated or deflected by the electric field of target nuclei. Because electrons can lose different fractions of their kinetic energy, bremsstrahlung forms a continuous X-ray spectrum from low energies up to the endpoint energy $E_{\max}$.

\subsection{Characteristic Radiation}

\textbf{Characteristic X-rays} are produced when an incident electron ejects an inner-shell electron from the target atom. An electron from a higher shell then falls into the vacancy and emits a photon with energy equal to the difference between shell binding energies:
\begin{equation}
    E_{\mathrm{char}} = E_{\mathrm{outer}} - E_{\mathrm{inner}}.
\end{equation}
These photons have discrete energies that are characteristic of the target material.

\subsection{Factors Affecting X-Ray Output}

The X-ray beam depends on several parameters:
\begin{itemize}
    \item \textbf{Tube voltage, kVp:} controls maximum photon energy and beam penetrability.
    \item \textbf{Tube current, mA:} controls the rate of electron flow and therefore photon quantity.
    \item \textbf{Exposure time:} combines with current as mAs, which controls total photon number.
    \item \textbf{Target material:} affects X-ray production efficiency and characteristic energies.
    \item \textbf{Filtration:} removes low-energy photons and hardens the beam.
    \item \textbf{Collimation:} restricts the beam to the useful field.
    \item \textbf{Anode angle:} affects focal spot size, heat loading, and the heel effect.
\end{itemize}

\section{Schematic Diagram of X-Ray Sources}

\subsection{Basic X-Ray Tube Production Chain}

\begin{lstlisting}
Filament current heats cathode
        |
        v
Thermionic emission releases electrons
        |
        v
Focusing cup directs electrons toward focal spot
        |
        v
High voltage accelerates electrons
        |
        v
Electrons strike metal target/anode
        |
        +--> Bremsstrahlung continuum
        |
        +--> Inner-shell ionization of target atoms
                |
                v
            Characteristic X-ray lines
        |
        v
Tube exit window
        |
        v
Inherent and added filtration
        |
        v
Collimator
        |
        v
Useful X-ray beam
\end{lstlisting}

\subsection{Coolidge Tube}

The Coolidge tube is a thermionic X-ray tube. It uses a heated cathode filament, a vacuum envelope, and a metal anode target.

\begin{lstlisting}
Cathode                 Vacuum envelope                  Anode target
[filament] ---- e- beam -------------------------------> / W target /
    |                                                     /
[focusing cup]                                           /
                                                        /
                     X-rays exit through tube window --> useful beam
\end{lstlisting}

\subsection{Rotating-Anode Tube}

Modern medical X-ray systems usually use a rotating-anode tube. The anode rotates so that heat is distributed over a larger circular focal track.

\begin{lstlisting}
Cathode --> focused electron beam --> angled rotating W-Re anode track
                                      | rotor inside envelope |
                                      | stator coils outside  |
                                      v
                           heat spread over circular track
                                      |
                                      v
                         X-rays exit through window
                                      |
                                      v
                                  collimator
\end{lstlisting}

\subsection{Synchrotron X-Ray Source}

A synchrotron source produces X-rays when relativistic electrons are bent by magnets or pass through insertion devices.

\begin{lstlisting}
Electron gun -> linac/booster -> storage ring
                                      |
                                      v
             bending magnet / wiggler / undulator
                                      |
                                      v
                            synchrotron X-rays
                                      |
                                      v
                              beamline optics
                                      |
                                      v
                              experimental station
\end{lstlisting}

\section{Detailed Structure of a Medical Radiography X-Ray Machine}

A modern medical radiography unit is a complete system, not just a tube. It includes the high-voltage generator, control console, tube housing, X-ray tube insert, filtration, collimator, patient support, detector, image-processing workstation, and safety systems.

\subsection{System Block Diagram}

\begin{lstlisting}
Control console
(kVp, mA, time, AEC selection)
        |
        v
High-voltage generator
        |
        v
Tube head / tube housing
        |
        +--> Cathode filament and focusing cup
        +--> Rotating anode, rotor, and stator
        +--> Oil cooling and lead-lined housing
        |
        v
Exit window and added filtration
        |
        v
Collimator and light field
        |
        v
Patient on table or wall stand
        |
        v
Anti-scatter grid / bucky
        |
        v
Digital detector
        |
        v
Image acquisition workstation

Safety features:
beam-on lamp, audible exposure-end signal,
filtration interlock, field-size logic, and shielding.
\end{lstlisting}

\subsection{Main Components and Their Functions}

\subsubsection{Cathode}

The cathode contains one or more tungsten filaments. When heated, the filament emits electrons by thermionic emission. A focusing cup surrounds the filament and shapes the electron cloud so that electrons strike a small focal spot on the anode.

\subsubsection{Anode}

The anode is the positive target. It receives the accelerated electrons and converts their kinetic energy mostly into heat and partly into X-rays. Tungsten is commonly used because it has high atomic number, high melting point, and good X-ray production efficiency. In a rotating-anode tube, the target track rotates to spread heat over a larger area.

\subsubsection{Vacuum Envelope}

The cathode and anode are sealed inside an evacuated glass or metal-ceramic envelope. The vacuum prevents electrons from colliding with gas molecules and allows the electron beam to travel efficiently from cathode to anode.

\subsubsection{Tube Housing}

The tube housing surrounds the tube insert. It provides mechanical support, electrical insulation, heat removal, and radiation shielding. Oil inside the housing transfers heat and insulates high-voltage parts. Lead lining reduces leakage radiation.

\subsubsection{High-Voltage Generator}

The high-voltage generator supplies the potential difference between cathode and anode. It also supplies filament power. Modern high-frequency generators provide a nearly constant tube voltage, improving beam consistency and efficiency.

\subsubsection{Filtration}

Filtration removes low-energy photons that would be absorbed near the patient's skin without contributing much to image formation. Aluminum is a common filter material in diagnostic radiography.

\subsubsection{Collimator}

The collimator uses lead shutters to limit the beam to the region of interest. A lamp and mirror project a visible light field that corresponds approximately to the X-ray field. Collimation reduces patient dose and improves image contrast by reducing scatter.

\subsubsection{Detector}

The detector records the X-ray pattern after the beam passes through the patient. Digital radiography systems use electronic detectors and image-processing workstations. The detector may be located in a table bucky or wall stand.

\subsection{Operation Sequence}

\begin{enumerate}[label=\arabic*.]
    \item The operator positions the patient and detector.
    \item The operator selects kVp, mA, exposure time or mAs, focal spot, and field size.
    \item The tube rotor is prepared and brought to speed.
    \item The cathode filament is heated.
    \item The high-voltage generator applies the selected tube potential.
    \item Electrons accelerate from cathode to anode.
    \item X-rays are produced at the anode by bremsstrahlung and characteristic processes.
    \item The beam passes through filtration and collimation.
    \item X-rays pass through the patient and are differentially attenuated.
    \item The detector records the transmitted pattern.
    \item The exposure ends by timer or automatic exposure control.
    \item The image is processed and displayed.
\end{enumerate}

\subsection{Typical Operating Parameters}

Representative medical radiography values include:
\begin{itemize}
    \item Tube voltage: approximately $40$--$150\,\mathrm{kVp}$.
    \item Photon energies: mostly in the tens of keV after filtration.
    \item Focal spots: often around $0.6\,\mathrm{mm}$ and $1.2\,\mathrm{mm}$ for dual-focus tubes.
    \item Filtration: commonly aluminum-equivalent filtration in general radiography.
    \item Exposure: controlled by mA, exposure time, and mAs.
\end{itemize}

\subsection{Cooling and Heat Management}

Most of the electron kinetic energy deposited in the anode becomes heat rather than X-rays. Therefore, heat management is the dominant engineering problem in X-ray tube design. A rotating anode distributes heat around a circular track. Oil in the housing removes heat and provides electrical insulation. Tube rating charts and cooling curves are used to prevent overheating.

\subsection{Safety Features}

Important safety features include:
\begin{itemize}
    \item lead shielding in the housing,
    \item collimation to restrict field size,
    \item filtration to remove low-energy photons,
    \item beam-on indicator lights,
    \item audible exposure termination signals,
    \item interlocks for filtration and field-size logic,
    \item protected operator barriers,
    \item quality control and calibration.
\end{itemize}

\section{Uses of X-Ray Machines}

X-ray machines are widely used in medicine, science, and industry. Important uses include:
\begin{enumerate}[label=\arabic*.]
    \item \textbf{Medical radiography:} imaging bones, chest, abdomen, trauma, and foreign bodies.
    \item \textbf{Dental imaging:} detecting cavities, tooth position, and jaw structure.
    \item \textbf{Fluoroscopy:} real-time imaging for procedures and contrast studies.
    \item \textbf{Computed tomography:} cross-sectional imaging using rotating X-ray sources and detectors.
    \item \textbf{Security scanning:} inspecting baggage and cargo.
    \item \textbf{Industrial radiography:} detecting cracks, weld defects, and internal flaws.
    \item \textbf{X-ray diffraction:} studying crystal structure.
    \item \textbf{X-ray fluorescence:} identifying elemental composition.
\end{enumerate}

\section{Shielding and Protection}

X-ray shielding depends on photon energy and material. High-$Z$ materials such as lead attenuate X-rays efficiently and are used in barriers, aprons, and tube housings. Concrete is useful for structural shielding in walls. Aluminum and copper are often used as beam filters rather than room-shielding materials.

For a narrow beam, shielding follows
\begin{equation}
    I = I_0 e^{-\mu x},
\end{equation}
where $I_0$ is incident intensity, $I$ is transmitted intensity, $\mu$ is the attenuation coefficient, and $x$ is absorber thickness. Using the half-value layer,
\begin{equation}
    I = I_0\left(\frac{1}{2}\right)^{x/\mathrm{HVL}}.
\end{equation}

\section{Comparison of X-Ray Source Types}

\begin{longtable}{p{0.18\textwidth} p{0.21\textwidth} p{0.19\textwidth} p{0.18\textwidth} p{0.16\textwidth}}
\caption{Comparison of common X-ray source types.}\\
\toprule
\textbf{Source type} & \textbf{Core target geometry} & \textbf{Spectrum character} & \textbf{Main strength} & \textbf{Typical use}\\
\midrule
\endfirsthead

\toprule
\textbf{Source type} & \textbf{Core target geometry} & \textbf{Spectrum character} & \textbf{Main strength} & \textbf{Typical use}\\
\midrule
\endhead

Classical Coolidge stationary-anode tube & Fixed metal target in evacuated thermionic tube & Bremsstrahlung continuum plus characteristic lines & Simple construction & Early radiography, low-duty applications, teaching\\

Rotating-anode diagnostic tube & Angled rotating target track driven by rotor and stator & Same physics as Coolidge tube, with greater heat capacity & Higher current and shorter exposures & General radiography, fluoroscopy, CT\\

Synchrotron source & Relativistic electrons in a storage ring bent by magnets or insertion devices & Highly bright, collimated, tunable radiation & Extremely high brightness and tunability & Crystallography, spectroscopy, micro-imaging, materials science\\
\bottomrule
\end{longtable}

\section{Electronic Configurations and Valencies}

The requested electronic configurations and common valencies are shown in Table~\ref{tab:configs}. Valency is written here in the ordinary chemistry sense of common combining capacity or common oxidation state.

\begin{table}[h!]
\centering
\caption{Electronic configurations and valencies of selected elements.}
\label{tab:configs}
\begin{tabularx}{\textwidth}{l c X l X}
\toprule
\textbf{Element} & \textbf{Atomic number} & \textbf{Electronic configuration} & \textbf{Common valency} & \textbf{Common oxidation states}\\
\midrule
Lead (Pb) & 82 & $[\mathrm{Xe}]\,4f^{14}5d^{10}6s^2 6p^2$ & 2 or 4 & $+2$, $+4$\\
Zinc (Zn) & 30 & $[\mathrm{Ar}]\,3d^{10}4s^2$ & 2 & $+2$\\
Titanium (Ti) & 22 & $[\mathrm{Ar}]\,3d^2 4s^2$ & 3 or 4 & $+2$, $+3$, $+4$; most common is $+4$\\
Aluminum (Al) & 13 & $[\mathrm{Ne}]\,3s^2 3p^1$ & 3 & $+3$\\
\bottomrule
\end{tabularx}
\end{table}

\section*{Conclusion}

Radiation includes electromagnetic waves and energetic particles. Its most important practical classification is whether it is ionizing or non-ionizing. X-rays are ionizing electromagnetic radiation with short wavelength, high frequency, and high photon energy. They are produced in X-ray tubes when electrons emitted from a heated cathode are accelerated into a metal anode, where they generate a continuous bremsstrahlung spectrum and discrete characteristic X-ray lines. A medical radiography system uses a rotating-anode tube, high-voltage generator, tube housing, filtration, collimation, detector, and safety controls to form useful diagnostic images while limiting unnecessary exposure.

\section*{References}

\begin{enumerate}[label={[\arabic*]}]
    \item Centers for Disease Control and Prevention, \textit{About Non-Ionizing Radiation}. \url{https://www.cdc.gov/radiation-health/about/non-ionizing-radiation.html}
    \item World Health Organization, \textit{Ionizing Radiation and Health Effects}. \url{https://www.who.int/news-room/fact-sheets/detail/ionizing-radiation-and-health-effects}
    \item NASA/GSFC, \textit{The Electromagnetic Spectrum}. \url{https://imagine.gsfc.nasa.gov/science/toolbox/emspectrum1.html}
    \item NASA/GSFC, \textit{Electromagnetic Spectrum Chart}. \url{https://imagine.gsfc.nasa.gov/science/toolbox/spectrum_chart.html}
    \item NASA HEASARC, \textit{The Electromagnetic Spectrum}. \url{https://heasarc.gsfc.nasa.gov/docs/heasarc/headates/spectrum.html}
    \item National Center for Biotechnology Information, \textit{X-ray Imaging---Medical Imaging Systems}. \url{https://www.ncbi.nlm.nih.gov/books/NBK546155/}
    \item National Center for Biotechnology Information, \textit{X-ray Production}. \url{https://www.ncbi.nlm.nih.gov/books/NBK564423/}
    \item National Institute of Standards and Technology, \textit{XCOM: Photon Cross Sections Database}. \url{https://www.nist.gov/pml/xcom-photon-cross-sections-database}
    \item National Institute of Standards and Technology, \textit{X-Ray Mass Attenuation Coefficients}. \url{https://www.nist.gov/pml/x-ray-mass-attenuation-coefficients}
    \item International Atomic Energy Agency, \textit{Radiation Protection of Staff}. \url{https://www.iaea.org/resources/rpop/radiation-protection-of-staff}
    \item U.S. Food and Drug Administration, \textit{Radiography}. \url{https://www.fda.gov/radiation-emitting-products/medical-x-ray-imaging/radiography}
    \item U.S. Food and Drug Administration, \textit{Performance Standard for Diagnostic X-Ray Systems and Their Major Components}. \url{https://www.fda.gov/media/78382/download}
    \item Stanford Synchrotron Radiation Lightsource, \textit{Hardware Source}. \url{https://www-ssrl.slac.stanford.edu/smb-saxs/content/hardware-source}
    \item Royal Society of Chemistry, \textit{Lead}. \url{https://periodic-table.rsc.org/element/82/lead}
    \item Royal Society of Chemistry, \textit{Zinc}. \url{https://periodic-table.rsc.org/element/30/zinc}
    \item Royal Society of Chemistry, \textit{Titanium}. \url{https://periodic-table.rsc.org/element/22/titanium}
    \item Royal Society of Chemistry, \textit{Aluminium}. \url{https://periodic-table.rsc.org/element/13/aluminium}
\end{enumerate}
\end{document}